% Table 10: Vol-managed portfolio metrics by regime

\begin{table}[htbp]
\centering
\caption{Volatility-Managed Portfolios: Risk-Adjusted Performance}
\label{tab:volmanaged}
\small
\setlength{\tabcolsep}{3pt}
\resizebox{\linewidth}{!}{%
\begin{tabular}{llccccc}
\toprule
Regime & Portfolio & Ann.\ Ret. & Ann.\ Vol. & Sharpe & Max DD & CER \\
\midrule
Full sample & Unmanaged (buy-and-hold) & +10.94\% & +14.89\% & +0.73 & -28.45\% & +5.40\% \\
 & A-managed & +9.13\% & +12.91\% & +0.71 & -19.89\% & +4.96\% \\
 & A1-managed & +9.22\% & +12.84\% & +0.72 & -18.56\% & +5.10\% \\
 & C-managed & +9.27\% & +12.72\% & +0.73 & -17.92\% & +5.23\% \\
\midrule
Low VIX (Q1) & Unmanaged (buy-and-hold) & +6.71\% & +7.69\% & +0.87 & -13.19\% & +5.23\% \\
 & A-managed & +9.29\% & +11.65\% & +0.80 & -20.54\% & +5.90\% \\
 & A1-managed & +10.65\% & +12.12\% & +0.88 & -20.18\% & +6.98\% \\
 & C-managed & +11.37\% & +12.41\% & +0.92 & -20.21\% & +7.52\% \\
\midrule
High VIX (Q4) & Unmanaged (buy-and-hold) & +21.02\% & +23.19\% & +0.91 & -24.12\% & +7.58\% \\
 & A-managed & +11.63\% & +13.55\% & +0.86 & -16.71\% & +7.04\% \\
 & A1-managed & +10.87\% & +12.39\% & +0.88 & -14.68\% & +7.04\% \\
 & C-managed & +11.23\% & +12.08\% & +0.93 & -14.35\% & +7.58\% \\
\midrule
COVID 2020 & Unmanaged (buy-and-hold) & +19.62\% & +30.07\% & +0.65 & -11.96\% & -2.98\% \\
 & A-managed & +11.03\% & +12.09\% & +0.91 & -4.22\% & +7.38\% \\
 & A1-managed & +11.48\% & +10.82\% & +1.06 & -4.33\% & +8.55\% \\
 & C-managed & +15.54\% & +11.39\% & +1.36 & -5.04\% & +12.30\% \\
\midrule
\bottomrule
\end{tabular}%
}
\begin{tablenotes}\small
\item Notes: Portfolios are equal-weighted across 115 stocks and rebalanced
every five trading days, in line with the forecast sampling interval.
The volatility-managed strategy uses stock-level weights
$w_{m,i,t}=c_{m,i}/\hat\sigma^2_{m,i,t}$, where
$\hat\sigma^2_{m,i,t}=\exp(\hat y_{m,i,t})$ is the exponentiated forecast
of log mean $RV^{PK}$ over $[t+1,t+5]$. The constant $c_{m,i}$ normalizes
the managed stock return so that $\mathrm{Var}(f_{m,i})=\mathrm{Var}(r_i)$.
Annualization assumes $252/5$ five-day periods per year. Certainty-equivalent
returns (CER) use a mean--variance investor with risk aversion $\gamma=5$.
The table reports descriptive performance statistics.
\end{tablenotes}
\end{table}